\chapter{METHODOLOGY}
\label{ch:methodology}

\section{Overview}

This chapter describes the methodology used to develop and evaluate deep learning models for brain tumor detection with explainability analysis. The pipeline consists of four main components: (1) data preprocessing and patient-centric loading, (2) model architecture and training with transfer learning, (3) comprehensive evaluation with F1-score metrics, and (4) explainability visualization using Gradient-weighted Class Activation Mapping (Grad-CAM).

\section{Dataset}

\rev{We use the BraTS 2021 Task 1 dataset containing 1,251 glioblastoma patients with multi-modal MRI sequences. Each patient scan includes four modalities: T1-weighted (T1), T1-weighted contrast-enhanced (T1ce), T2-weighted (T2), and Fluid Attenuated Inversion Recovery (FLAIR). All images are provided at 240×240×155 voxels with 1mm³ isotropic resolution, along with expert-annotated tumor segmentation masks.

\textbf{Modality Selection:} We use three MRI modalities (T1ce, T2, FLAIR) as input channels, which directly matches the 3-channel input expected by ImageNet-pretrained models. This allows us to use standard pretrained weights without modification. The T1 (non-contrast) modality is excluded as T1ce provides similar anatomical information with additional contrast enhancement useful for tumor visualization.

\textbf{Data Splitting:} We implement patient-level splitting with a fixed random seed (42) to ensure reproducibility and prevent data leakage. Patient-level splitting is critical---all slices from a single patient remain in the same split to ensure reliable generalization estimates and avoid pseudo-replication bias. Table \ref{tab:data_split} shows the distribution of patients and slices across splits.

\begin{table}[htbp]
\centering
\caption{Dataset distribution with tumor/non-tumor slice counts}
\label{tab:data_split}
\begin{tabular}{lcccc}
\toprule
\textbf{Split} & \textbf{Patients} & \textbf{Total Slices} & \textbf{Tumor Slices} & \textbf{Non-Tumor Slices} \\
\midrule
Training & 877 (70\%) & 135,935 & 56,672 (41.69\%) & 79,263 (58.31\%) \\
Validation & 187 (15\%) & 28,985 & 12,017 (41.46\%) & 16,968 (58.54\%) \\
Test & 187 (15\%) & 28,985 & 12,748 (43.98\%) & 16,237 (56.02\%) \\
\midrule
\textbf{Total} & \textbf{1,251} & \textbf{193,905} & \textbf{81,437} & \textbf{112,468} \\
\bottomrule
\end{tabular}
\end{table}}

\section{Data Preprocessing and Loading}

\rev{Unlike previous work focused on MGMT methylation prediction or segmentation, our task is \textbf{slice-level binary tumor detection}: classifying individual 2D axial slices as either containing tumor tissue (positive) or tumor-free (negative). This formulation provides abundant training data ($\sim$155 slices × 1,251 patients $\approx$ 194,000 samples), enables direct comparison of CNN vs Transformer attention mechanisms, and facilitates explainability analysis through 2D heatmap visualization.

\textbf{Slice Labeling:} Ground truth labels are derived from expert-annotated segmentation masks. For each 2D axial slice, if any pixel in the corresponding segmentation mask has a non-zero value (indicating tumor tissue), the slice is labeled as positive (tumor present, label = 1). Otherwise, the slice is labeled as negative (tumor-free, label = 0).

\textbf{Patient-Centric Batch Loading:} To handle the large dataset efficiently while maintaining patient-level splits, we implement a patient-centric batching strategy where each training step loads all $\sim$155 slices from a single patient. This approach balances computational efficiency with the need to process all patient slices for comprehensive tumor detection evaluation.

\textbf{Preprocessing Pipeline:} Each 2D slice undergoes: (1) modality stacking---extracting corresponding slices from T1ce, T2, and FLAIR volumes and stacking into a 3-channel image; (2) Z-score normalization per channel using foreground voxels: $x_{\text{norm}} = \frac{x - \mu}{\sigma + \epsilon}$; (3) resizing from 240×240 to 224×224 using bilinear interpolation; and (4) conversion to PyTorch tensors with channel-first format (C×H×W). No data augmentation is applied during training to maintain consistent evaluation of model architectures.}

\section{Model Architectures}

\rev{We compare three state-of-the-art architectures representing different design philosophies for image classification. Each architecture was selected for its unique characteristics and significance in deep learning research:

\subsection{ResNet-50: Deep Residual CNN}

ResNet-50 (23.6M parameters) introduced skip connections that revolutionized deep learning by enabling training of very deep networks \cite{he2016deep}. \textbf{Significance:} ResNet's residual learning framework solved the degradation problem in deep networks, achieving state-of-the-art results on ImageNet and becoming the de facto baseline for medical imaging transfer learning.

Architecture features:
\begin{itemize}
    \item \textbf{Backbone}: 50-layer CNN with residual blocks (conv1 + 4 stages of bottleneck blocks)
    \item \textbf{Skip connections}: Identity mappings allowing gradients to flow directly through the network
    \item \textbf{Pretrained weights}: ImageNet-1K (1.2M images, 1000 classes)
    \item \textbf{Classification head}: Linear(2048 → 1) with sigmoid activation
    \item \textbf{Target layer for Grad-CAM}: \texttt{layer4.2.conv3} (last convolutional layer)
\end{itemize}

\subsection{EfficientNet-B2: Compound-Scaled CNN}

EfficientNet-B2 (7.7M parameters) uses compound scaling to systematically balance network depth, width, and resolution \cite{tan2019efficientnet}. \textbf{Significance:} EfficientNet demonstrated that balanced scaling outperforms single-dimension scaling, achieving higher accuracy with fewer parameters than ResNet-50, making it particularly valuable for resource-constrained medical imaging applications.

Architecture features:
\begin{itemize}
    \item \textbf{Backbone}: Mobile inverted bottleneck convolution (MBConv) blocks with squeeze-and-excitation
    \item \textbf{Compound scaling}: Depth × 1.1, width × 1.1, resolution × 1.2 relative to B0
    \item \textbf{Pretrained weights}: ImageNet-1K with AutoML-discovered architecture
    \item \textbf{Classification head}: Linear(1408 → 1) with sigmoid activation
    \item \textbf{Target layer for Grad-CAM}: \texttt{features.8.0} (final MBConv block)
\end{itemize}

\subsection{DeiT-Base: Vision Transformer}

DeiT-Base (Data-efficient Image Transformer, 86M parameters) applies self-attention mechanisms from NLP to images \cite{touvron2021training}. \textbf{Significance:} DeiT proved that transformers can match CNN performance on ImageNet using only standard training data (without the 300M images required by original ViT), enabling exploration of fundamentally different attention mechanisms for medical imaging.

Architecture features:
\begin{itemize}
    \item \textbf{Backbone}: 12-layer transformer encoder with 12 attention heads
    \item \textbf{Input}: 224×224×3 image split into 16×16 patches (196 patches total)
    \item \textbf{Pretrained weights}: ImageNet-1K with knowledge distillation training
    \item \textbf{Classification head}: Linear(768 → 1) with sigmoid activation
    \item \textbf{Target layer for Grad-CAM}: \texttt{patch\_embed.proj} (patch embedding convolution)
\end{itemize}

\begin{figure}[htbp]
\centering
\includegraphics[width=\textwidth]{architectures/vit_arch.png}
\caption{Vision Transformer architecture. The image is split into fixed-size patches, linearly embedded with position encodings, and processed by a standard Transformer encoder. A learnable classification token aggregates information for the final prediction. Figure inspired by Dosovitskiy et al. (2021).}
\label{fig:vit_arch}
\end{figure}

All three models use standard 3-channel input (T1ce, T2, FLAIR), matching the input format of ImageNet-pretrained models.}

\section{Training Protocol}

\subsection{Two-Stage Transfer Learning}

We employ a two-stage training strategy to leverage pretrained features effectively:

\begin{itemize}
    \item \textbf{Stage 1 (Epochs 0-4)}: Backbone frozen, only classification head trainable
    \item \textbf{Stage 2 (Epochs 5-9)}: Full network fine-tuning with unfrozen backbone
    \item \textbf{Total epochs}: 10 (sufficient for convergence on large dataset)
\end{itemize}

\subsection{Optimization Configuration}

\begin{itemize}
    \item \textbf{Optimizer}: AdamW with $\beta_1=0.9$, $\beta_2=0.999$
    \item \textbf{Learning rate}: Initial $10^{-3}$, decayed with Cosine Annealing schedule ($T_{\max}=10$ epochs, $\eta_{\min}=10^{-5}$)
    \item \textbf{Weight decay}: $10^{-4}$ (L2 regularization)
    \item \textbf{Loss function}: Binary Cross-Entropy with Logits (BCEWithLogitsLoss)
    \item \textbf{Batch size}: 1 patient ($\sim$155 slices per batch)
    \item \textbf{Gradient accumulation}: None (sufficient with patient-centric batching)
\end{itemize}

\subsection{Checkpointing and Early Stopping}

Models are checkpointed based on validation accuracy (not loss) to ensure best generalization:
\begin{itemize}
    \item \textbf{Checkpoint criterion}: Maximum validation accuracy
    \item \textbf{Save strategy}: Best + last checkpoint
    \item \textbf{Test evaluation}: Using best validation checkpoint only
\end{itemize}

\section{Evaluation Metrics}

To comprehensively assess model performance on the slice-level tumor detection task, we track the following metrics:

\subsection{Primary Metrics}

\begin{itemize}
    \item \textbf{Accuracy}: Overall classification correctness $= \frac{TP + TN}{TP + TN + FP + FN}$
    \item \textbf{F1-Score}: Harmonic mean of precision and recall $= \frac{2 \cdot \text{Precision} \cdot \text{Recall}}{\text{Precision} + \text{Recall}}$
\end{itemize}

where Precision $= \frac{TP}{TP + FP}$ (proportion of positive predictions that are correct) and Recall $= \frac{TP}{TP + FN}$ (proportion of actual positives that are detected). TP, TN, FP, FN denote true positives, true negatives, false positives, and false negatives respectively.

F1-score is particularly important for imbalanced datasets, providing a balanced assessment of model performance across positive and negative classes.

\subsection{Metrics Tracked Across Training}

All metrics are computed and logged for three phases:
\begin{itemize}
    \item \textbf{Training}: Computed on training batches (logged every epoch)
    \item \textbf{Validation}: Computed on full validation set (logged every epoch)
    \item \textbf{Test}: Computed once on held-out test set using best checkpoint
\end{itemize}

This comprehensive tracking enables analysis of overfitting, convergence behavior, and final generalization performance.

\section{Explainability: Grad-CAM Implementation}

To provide interpretable predictions and analyze where models focus attention, we implement Gradient-weighted Class Activation Mapping (Grad-CAM) \cite{selvaraju2017grad}.

\subsection{Grad-CAM Algorithm}

Grad-CAM generates a coarse localization heatmap highlighting image regions most important for the predicted class. The algorithm computes:

\begin{equation}
L^c_{\text{Grad-CAM}} = \text{ReLU}\left(\sum_k \alpha_k^c \mathbf{A}^k\right)
\end{equation}

where $\alpha_k^c = \frac{1}{HW} \sum_i \sum_j \frac{\partial y^c}{\partial A^k_{ij}}$ represents the importance of feature map $k$ for class $c$, and $\mathbf{A}^k$ is the $k$-th activation map from the target convolutional layer.

\subsection{Implementation Details}

\textbf{Target layer selection:}
\begin{itemize}
    \item \textbf{EfficientNet-B2}: \texttt{features.8.0} (final MBConv block)
    \item \textbf{ResNet-50}: \texttt{layer4.2.conv3} (last convolutional layer)
    \item \textbf{DeiT-Base}: \texttt{patch\_embed.proj} (patch embedding convolution)
\end{itemize}

\textbf{Heatmap generation pipeline:}
\begin{enumerate}
    \item Forward pass with gradient tracking on target layer activations
    \item Backward pass to compute gradients w.r.t. positive class prediction
    \item Global average pooling of gradients to obtain importance weights
    \item Weighted sum of activation maps, followed by ReLU (removes negative activations)
    \item Bilinear upsampling to original input resolution (224×224)
    \item Gaussian smoothing ($\sigma=3$) to reduce blockiness from low-resolution activations
\end{enumerate}

\subsection{Visualization Strategy}

For each sample, we generate a 3-column visualization:

\begin{enumerate}
    \item \textbf{Column 1 - Ground Truth}: MRI slice (T1ce) with green tumor segmentation overlay (alpha=0.7)
    \item \textbf{Column 2 - GradCAM Attention}: MRI slice with red heatmap overlay showing model attention (alpha=0.7, hot colormap)
    \item \textbf{Column 3 - Combined}: MRI slice with both green segmentation and red attention heatmap
\end{enumerate}

\textbf{Robust normalization}: To ensure brain tissue remains visible even with bright artifacts, we use \textbf{99.5th percentile normalization}:
\begin{equation}
I_{\text{norm}} = \text{clip}\left(\frac{I}{\text{percentile}_{99.5}(I) + \epsilon}, 0, 1\right)
\end{equation}

This approach prevents outlier pixels from compressing the dynamic range, ensuring both brain anatomy and heatmaps are visible.

\subsection{Sample Selection for Analysis}

From each model's test set predictions, we select high-quality visualizations using the following criteria:

\begin{itemize}
    \item \textbf{True positives only}: Focus on correct detections to understand what models learn
    \item \textbf{Substantial tumor size}: Filter slices with tumor area $<$ 50 pixels (reduces noise from tiny tumors)
    \item \textbf{Non-empty slices}: Exclude all-black or nearly empty slices (max pixel intensity $<$ 10)
    \item \textbf{Diverse patient coverage}: Select 2 slices from 20 different patients per model (40 samples total)
\end{itemize}

Note that this filtering applies only to visualization sample selection, not to model training. During training, the model learns from all tumor slices regardless of size to ensure robustness to varying tumor appearances. The filtering here ensures that Grad-CAM visualizations show clear, representative examples for qualitative analysis.

This selection ensures that visualizations represent meaningful model behavior on substantial, well-defined tumor regions across diverse patient cases.

\section{Reproducibility}

To ensure reproducibility, we:

\begin{itemize}
    \item \textbf{Fix random seeds}: PyTorch, NumPy, and Python random (seed=42)
    \item \textbf{Deterministic algorithms}: Enable deterministic CUDA operations where possible
    \item \textbf{Version control}: All code tracked in Git repository
    \item \textbf{Hyperparameter logging}: Training configuration saved with each experiment
    \item \textbf{Checkpoint saving}: Model weights, optimizer state, and metrics preserved
    \item \textbf{Environment specification}: requirements.txt with exact package versions
\end{itemize}

All experiments use identical data splits, preprocessing, and training protocols to enable fair comparison between ResNet50 and EfficientNet-B2.

\section{Summary}

This chapter presented a comprehensive methodology for slice-level brain tumor detection with explainability, featuring:
\begin{enumerate}
    \item Rigorous data handling with patient-level splitting to avoid pseudo-replication
    \item Three modern architectures: EfficientNet-B2 (compound scaling CNN), ResNet-50 (residual CNN), and DeiT-Base (vision transformer)
    \item Rationale for 3-channel RGB approach (T1ce, T2, FLAIR) compatible with ImageNet-pretrained models
    \item Systematic two-stage training protocol with comprehensive evaluation metrics (accuracy, precision, recall, F1)
    \item Integration of Grad-CAM visualization with percentile-based normalization for clinical interpretability
    \item Reproducible implementation with fixed seeds and version control
\end{enumerate}

The next chapter presents the experimental results obtained using this methodology.
